\SourceLeaf{272}
\Article{5}
\indent I accompany the instructions in this whole memoir with a numerical example worked out with the greatest precision. For others who wish to use the method presented here to process their measurements, it can serve either as a model for computing the necessary auxiliary tables or directly as an aid for a large part of the temperate zone. In most cases, however, \textit{much} less precision will suffice.

I choose \(52^0\ 40'\) as the standard latitude, roughly the middle parallel of the Kingdom of Hanover. In some respects it is more advantageous to choose a round number for the standard latitude on the sphere than for that on the ellipsoid, so I set
\[
Q=52^0\ 40'\ 0''
\]
\indent I carry out the calculation using the latest dimensions of the Earth that \textsc{Bessel} derived from arc measurements (Astronomische Nachrichten, vol. 19, p. 116). Taking the toise as the unit gives
\[
\begin{aligned}
\log a&=6{,}5148235337\\
\log\cos\varphi&=9{,}9985458202
\end{aligned}
\]
With the help of ten-digit logarithms we obtain
\[
\begin{aligned}
\varphi&=4^0\ 41'\ 9''98262\\
\log e&=8{,}9122052079\\
\zeta&=1^0\ 43'\ 26''80402\\
\eta&=2\quad15\quad42\ 34083\\
P&=52\quad42\quad2{,}53251\\
\log\alpha&=0{,}0001966553\\
\theta&=3^0\ 43'\ 34''24669\\
\log\frac{1}{k}&=0{,}0016708804\\
\log A&=6{,}5152074703
\end{aligned}
\]
Taking the French legal metre as the unit gives
\[
\log A=6{,}8050274003
\]
